\documentclass[11pt]{article}

% ===================== PACKAGES =====================
\usepackage[a4paper,margin=1in]{geometry}
\usepackage{amsmath,amssymb}
\usepackage{enumitem}
\usepackage{fancyhdr}
\usepackage{xcolor} 

% ===================== HEADER & FOOTER =====================
\pagestyle{fancy}
\fancyhf{}
\lhead{CSE 400: Fundamentals of Probability in Computing}
\rhead{Homework Assignment-2}
\cfoot{\thepage}

% ===================== CUSTOM COMMANDS =====================
\newcommand{\solution}{
    \vspace{0.5cm}
    \noindent\textbf{\textcolor{blue}{Solution:}}
    \vspace{0.2cm}
}

% ===================== TITLE =====================
\title{
    \normalsize School of Engineering and Applied Science (SEAS), Ahmedabad University \\
    \vspace{0.2cm}
    \textbf{CSE 400: Fundamentals of Probability in Computing}\\
    \Large Lecture scribe submission
}
\author{}
\date{}

\begin{document}
\maketitle

\vspace{-2cm}
\begin{center}
    \begin{tabular}{ll}
        \textbf{Name:} & Ishika Patel \\[1.5ex]
        \textbf{Enrollment No:} & AU2440088 \\[1.5ex]
        \textbf{Group:} & s1\_g1\_fin \\[1.5ex]
        \textbf{Email:} & ishika.p3@ahduni.edu.in \\[1.5ex]
        \textbf{Lecture:} & 6 \\[1.5ex]
        \textbf{Date of Submission:} & February 7\textsuperscript{th}, 2026 (11:59 PM)
    \end{tabular}
\end{center}

\hrule
\vspace{0.5cm}

\section{Random Variables}

\subsection{Definition}

A random variable $X$ on a sample space $\Omega$ is a function
\[
X : \Omega \rightarrow \mathbb{R}
\]
that assigns to each sample point $\omega \in \Omega$ a real number.

Until further notice, attention is restricted to discrete random variables, i.e., random variables that take values in a finite or countably infinite set.

Although $X$ maps $\Omega$ to $\mathbb{R}$, the actual set of values
\[
\{ X(\omega) : \omega \in \Omega \}
\]
is a discrete subset of $\mathbb{R}$.

\subsection{Distribution of a Random Variable}

The distribution of a random variable can be visualized as a bar diagram:
\begin{itemize}
    \item The x-axis represents the values the random variable can take.
    \item The height of the bar at value $a$ is $\Pr[X = a]$.
\end{itemize}

Each probability is computed from the probability of the corresponding event in the sample space.

\section{Guide to Selecting a Probability Distribution}

\subsection{Discrete Random Variable}
\begin{itemize}
    \item Countable support
    \item Probability Mass Function
    \item Probabilities assigned to single values
    \item Each possible value has strictly positive probability
\end{itemize}

\subsection{Continuous Random Variable}
\begin{itemize}
    \item Uncountable support
    \item Probability Density Function
    \item Probabilities assigned to intervals of values
    \item Each possible value has zero probability
\end{itemize}

\section{Examples of Random Variables}

\subsection{Example 1: Tossing 3 Fair Coins}

Let $Y$ denote the number of heads that appear when 3 fair coins are tossed.

Possible values:
\[
Y \in \{0,1,2,3\}
\]

Probabilities:
\[
\Pr(Y=0) = \Pr(t,t,t)
\]
\[
\Pr(Y=1) = \Pr\{(t,t,h),(t,h,t),(h,t,t)\}
\]
\[
\Pr(Y=2) = \Pr\{(h,h,t),(h,t,h),(t,h,h)\}
\]
\[
\Pr(Y=3) = \Pr(h,h,h)
\]

Since $Y$ must take one of the values 0 through 3:
\[
\sum_{i} \Pr(Y=i) = 1
\]

\section{Probability Mass Function (PMF)}

\subsection{Definition}

A random variable that can take on at most a countable number of possible values is said to be discrete.

Let $X$ be a discrete random variable with range
\[
R_X = \{x_1, x_2, x_3, \ldots\}
\]
(finite or countably infinite).

The function
\[
p_X(x_k) = \Pr(X = x_k), \quad k = 1,2,3,\ldots
\]
is called the Probability Mass Function (PMF) of $X$.

Since $X$ must take one of these values:
\[
\sum_{k=1}^{\infty} p_X(x_k) = 1
\]

\subsection{PMF Example}

Given:
\[
p(i) = c e^{-i}, \quad i = 0,1,2,\ldots
\]

Since
\[
\sum_{i=0}^{\infty} p(i) = 1
\]
we have:
\[
c \sum_{i=0}^{\infty} e^{-i} = 1
\]

Using
\[
\sum_{i=0}^{\infty} e^{-i} = \frac{1}{1-e^{-1}}
\]
implies:
\[
c = e^{-1}
\]

Results:
\[
\Pr(X=0) = e^{-1}
\]
\[
\Pr(X>2) = 1 - \Pr(X \le 2) = 1 - \Pr(X=0) - \Pr(X=1) - \Pr(X=2)
\]

\section{Bayes' Theorem (Recap)}

\subsection{Bayes Formula (Proposition 3.1)}

\[
\Pr(A \mid B_i) =
\frac{\Pr(B_i \mid A)\Pr(A)}
{\sum_{i=1}^{n} \Pr(A \mid B_i)\Pr(B_i)}
\]

\[
\Pr(B_i): \text{ A priori probability}
\]
\[
\Pr(B_i \mid A): \text{ Posteriori probability}
\]

\subsection{Example 1: Auditorium with 30 Rows}

\begin{itemize}
    \item 30 rows
    \item Row 1 has 11 seats, Row 30 has 40 seats
    \item A row is selected uniformly
    \item A seat within the selected row is selected uniformly
\end{itemize}

Tasks:
\[
\Pr(\text{Seat 15} \mid \text{Row 20})
\]
\[
\Pr(\text{Row 20} \mid \text{Seat 15})
\]

\subsection{Example 2: Communication System (Receiver)}

Conditional probabilities:
\[
\Pr(0_r \mid 0_t) = 0.95, \quad \Pr(1_r \mid 0_t) = 0.05
\]
\[
\Pr(0_r \mid 1_t) = 0.10, \quad \Pr(1_r \mid 1_t) = 0.90
\]

Tasks:
\begin{itemize}
    \item Find $\Pr(0 \text{ received})$ and $\Pr(1 \text{ received})$
    \item Find posterior probabilities of transmitted bit given received bit
    \item Find overall probability of error
\end{itemize}

\section{Independent Events}

\subsection{Definition}

Two events $A$ and $B$ are independent if:
\[
\Pr(A \mid B) = \Pr(A), \quad \Pr(B \mid A) = \Pr(B)
\]

Hence:
\[
\Pr(A,B) = \Pr(A)\Pr(B)
\]

\subsection{Mutual Independence (Three Events)}

Events $A,B,C$ are mutually independent if:
\[
\Pr(A,B) = \Pr(A)\Pr(B)
\]
\[
\Pr(A,C) = \Pr(A)\Pr(C)
\]
\[
\Pr(B,C) = \Pr(B)\Pr(C)
\]
\[
\Pr(A,B,C) = \Pr(A)\Pr(B)\Pr(C)
\]

\section{Types of Discrete Random Variables}

\subsection{Bernoulli Random Variable}

Experiment outcome: Success or Failure

\[
X =
\begin{cases}
1 & \text{Success} \\
0 & \text{Failure}
\end{cases}
\]

PMF:
\[
p_X(0) = 1-p, \quad p_X(1) = p, \quad p \in (0,1)
\]

\subsection{Binomial Random Variable}

Experiment:
\begin{itemize}
    \item $n$ independent trials
    \item Success probability $p$
\end{itemize}

Random variable:
\[
X = \text{number of successes in } n \text{ trials}
\]

Notation:
\[
X \sim B(n,p)
\]

PMF:
\[
p(i) = \binom{n}{i} p^i (1-p)^{n-i}, \quad i = 0,1,\ldots,n
\]

\subsection{Geometric Random Variable}

Experiment:
\begin{itemize}
    \item Independent trials
    \item Success probability $p$
    \item Trials continue until first success
\end{itemize}

Random variable:
\[
X = \text{number of trials until first success}
\]

PMF:
\[
\Pr(X=n) = (1-p)^{n-1}p, \quad n = 1,2,\ldots
\]

\subsection{Example: Urn Problem}

\[
N \text{ white balls}, \quad M \text{ black balls}
\]

Replacement after each draw

\[
p = \frac{M}{M+N}
\]

\[
\Pr(X=n) = \left(\frac{N}{M+N}\right)^{n-1} \frac{M}{M+N}
\]

\subsection{Poisson Random Variable}

Definition: A random variable $X$ taking values $0,1,2,\ldots$ is Poisson with parameter $\lambda > 0$ if:
\[
\Pr(X=i) = \frac{e^{-\lambda}\lambda^i}{i!}, \quad i = 0,1,2,\ldots
\]

Normalization:
\[
\sum_{i=0}^{\infty} \Pr(X=i) = 1
\]

Note: Poisson RV approximates $B(n,p)$ when $n$ is large, $p$ is small, and $np$ is moderate.

\bigskip
\hrule
\vspace{0.2cm}
\begin{center}
    \small \textit{End of Submission}
\end{center}

\end{document}